% !TeX root = ../../金融经济学_熊和平_习题.tex
% 源：微信图片_20260928151601_3596_382.jpg（教材印张 331 起；性质 2 在扫描件末处截断）

\chapter{布朗运动及其简单描述}

\begin{tcolorbox}[summary,title=附录提要]
通常假设股票收益率由两部分确定：确定性的部分和随机性的部分，随机性由维纳过程或布朗运动刻画。参照 John Hull 的描述，称随机过程 $z$ 服从维纳过程，若 $z$ 具有如下两个性质。

\textbf{性质 1}：在很短的时间间隔 $\Delta t$ 内的增量 $\Delta z$ 可以表示为
\[
    \Delta z=\epsilon\sqrt{\Delta t}
    \tag{A.1}
\]
其中 $\epsilon$ 服从标准正态分布，即 $\epsilon\sim N(0,1)$；取 $\Delta t\to 0$ 时，(A.1) 可记为
\[
    \mathrm{d}z=\epsilon\sqrt{\mathrm{d}t}.
    \tag{A.2}
\]

\textbf{性质 2}：$z$ 在任何两个不相互重叠的时间区间上的增量相互独立。% 原书后续文字在扫描件末处截断，待补扫。
\end{tcolorbox}
